    \subsection{発展パターン・特殊性質型の演習}\label{節：発展パターン・特殊性質型の演習}
      ここでは,\sectionref{節：発展パターン}と\sectionref{節：特殊性質型}で扱った\bookterm{recurrence}{漸化式}を,\,15題で練習する.
      \bookterm{floor}{ガウス記号}では整数性や余り,多項間の関係では\bookterm{difference}{階差}や\bookterm{characteristic}{特性方程式},
      \bookterm{nonlinear}{非線形}な式では対称式や\bookterm{trigonometric}{三角関数}の恒等式が手掛かりになる.
      まず式のどの部分に着目したかを述べ,置換や場合分けを選ぼう.
      \bookterm{general}{一般項}の候補が見つかったら,\bookterm{initial}{初期条件}と\bookterm{recurrence}{漸化式}をともに満たすことを確認する.
      分数式では分母,\bookterm{trigonometric}{三角関数}では角度の定義域も確かめてほしい.
      問14・問15では初項を動かし,定\bookterm{sequence}{数列}・周期列・有界な\bookterm{sequence}{数列}が現れる条件を調べる.
      以下,特に指定しない限り,\,$\displaystyle n$は正の整数とする.
      $\displaystyle \lfloor x\rfloor$は$\displaystyle x$を超えない最大の整数を表す.
      周期の最小値とは,すべての添字で$\displaystyle a_{n+p}=a_n$となる正の整数$\displaystyle p$の最小値をいう.

      \subsubsection{問題}\label{節：発展パターン・特殊性質型の演習：問題}
        \begin{enumerate}[label=\textbf{問\arabic*.},leftmargin=3.8em,format=\bookitemlink{practice-special}{problem}{answer}]
          % E1: ガウス型・整数の半減
          \item \bookterm{sequence}{数列}$\displaystyle \{a_n\}$を
                \[
                  a_1=2026,\qquad
                  a_{n+1}=\left\lfloor\frac{a_n}{2}\right\rfloor
                \]
                によって定める.
                \par\smallskip\noindent
                (1) \bookterm{general}{一般項}を求めよ.
                \par\smallskip\noindent
                (2) 初めて$\displaystyle a_n=0$となる添字$\displaystyle n$を求めよ.
          % E2: ガウス型・三つずつの群
          \item \bookterm{sequence}{数列}$\displaystyle \{a_n\}$を
                \[
                  a_0=0,\qquad
                  a_{n+1}=a_n+\left\lfloor\frac{n+1}{3}\right\rfloor
                  \quad(n\ge0)
                \]
                によって定める.\,
                $\displaystyle n=3q+r$（$\displaystyle q\ge0$,\,$\displaystyle r=0,1,2$）とおき,
                \bookterm{general}{一般項}を$\displaystyle q,r$で表せ.
          % E3: ガウス型・余りの有限個の状態
          \item \bookterm{sequence}{数列}$\displaystyle \{a_n\}$を
                \[
                  a_1=1,\qquad
                  a_{n+1}=2a_n+1
                    -5\left\lfloor\frac{2a_n+1}{5}\right\rfloor
                \]
                によって定める.
                \par\smallskip\noindent
                (1) すべての項が$\displaystyle 0,1,2,3,4$のいずれかになることを示せ.
                \par\smallskip\noindent
                (2) \bookterm{general}{一般項}を,\,$\displaystyle n$を$\displaystyle 4$で割った余りで場合分けして表せ.
                また,周期の最小値を求めよ.
          % E4: 隣接4項間・三階差が定数
          \item \bookterm{sequence}{数列}$\displaystyle \{a_n\}$を
                \[
                  a_1=2,\qquad a_2=9,\qquad a_3=28,
                \]
                \[
                  a_{n+3}=3a_{n+2}-3a_{n+1}+a_n+6
                \]
                によって定める.
                \par\smallskip\noindent
                (1) $\displaystyle d_n=a_{n+2}-2a_{n+1}+a_n$とおき,\,$\displaystyle d_n$を求めよ.
                \par\smallskip\noindent
                (2) \bookterm{general}{一般項}を求めよ.
          % E5: 隣接4項間・異なる三つの特性根
          \item \bookterm{sequence}{数列}$\displaystyle \{a_n\}$を
                \[
                  a_1=3,\qquad a_2=6,\qquad a_3=14,
                \]
                \[
                  a_{n+3}=6a_{n+2}-11a_{n+1}+6a_n
                \]
                によって定める.
                \par\smallskip\noindent
                (1) $\displaystyle d_n=a_{n+2}-3a_{n+1}+2a_n$とおき,\,$\displaystyle d_n$を求めよ.
                \par\smallskip\noindent
                (2) \bookterm{general}{一般項}を求めよ.
          % E6: 虚数解型・周期6
          \item \bookterm{sequence}{数列}$\displaystyle \{a_n\}$を
                \[
                  a_1=0,\qquad a_2=1,\qquad
                  a_{n+2}=a_{n+1}-a_n
                \]
                によって定める.
                \par\smallskip\noindent
                (1) \bookterm{trigonometric}{三角関数}を用いて\bookterm{general}{一般項}を表せ.
                \par\smallskip\noindent
                (2) 周期の最小値を求めよ.
          % E7: 虚数解型・振幅が減衰する数列
          \item \bookterm{sequence}{数列}$\displaystyle \{a_n\}$を
                \[
                  a_1=1,\qquad a_2=0,\qquad
                  a_{n+2}=a_{n+1}-\frac12a_n
                \]
                によって定める.\bookterm{trigonometric}{三角関数}を用いて\bookterm{general}{一般項}を表し,
                $\displaystyle \lim_{n\to\infty}a_n$を求めよ.
          % E8: 非隣接2項間・奇数列と偶数列の逆数
          \item \bookterm{sequence}{数列}$\displaystyle \{a_n\}$を
                \[
                  a_1=1,\qquad a_2=2,\qquad
                  a_{n+2}=\frac{a_n}{1+a_n}
                \]
                によって定める.
                \par\smallskip\noindent
                (1) すべての項が正であることを示し,
                $\displaystyle a_{2k-1},a_{2k}$を$\displaystyle k$で表せ.
                ただし,\,$\displaystyle k$は正の整数とする.
                \par\smallskip\noindent
                (2) $\displaystyle \lim_{n\to\infty}a_n$を求めよ.
          % E9: 基本対称式型・平行移動して恒等式を使う
          \item \bookterm{sequence}{数列}$\displaystyle \{a_n\}$を
                \[
                  a_1=\frac72,\qquad
                  a_{n+1}=a_n^2-2a_n
                \]
                によって定める.
                \par\smallskip\noindent
                (1) $\displaystyle b_n=a_n-1$とおき,\,$\displaystyle b_n$の\bookterm{recurrence}{漸化式}を求めよ.
                \par\smallskip\noindent
                (2) 恒等式
                \[
                  \left(u+u^{-1}\right)^2-2=u^2+u^{-2}
                  \quad(u\ne0)
                \]
                を用いて$\displaystyle a_n$の\bookterm{general}{一般項}を求めよ.
          % E10: 基本対称式型・三乗の対称式
          \item \bookterm{sequence}{数列}$\displaystyle \{a_n\}$を
                \[
                  a_1=\frac52,\qquad
                  a_{n+1}=a_n^3-3a_n
                \]
                によって定める.
                \par\smallskip\noindent
                (1) 恒等式
                \[
                  \left(u+u^{-1}\right)^3
                  -3\left(u+u^{-1}\right)=u^3+u^{-3}
                \]
                を示せ.ただし,\,$\displaystyle u\ne0$とする.
                \par\smallskip\noindent
                (2) \bookterm{general}{一般項}を求めよ.
          % E11: 三角関数型・余弦の倍角と周期3
          \item \bookterm{sequence}{数列}$\displaystyle \{a_n\}$を
                \[
                  a_1=\cos\frac{2\pi}{7},\qquad
                  a_{n+1}=2a_n^2-1
                \]
                によって定める.
                \par\smallskip\noindent
                (1) \bookterm{general}{一般項}を求めよ.
                \par\smallskip\noindent
                (2) 周期の最小値を求めよ.
          % E12: 三角関数型・正接の定義域と周期4
          \item 初項と次の項を計算する規則を
                \[
                  a_1=\tan\frac{\pi}{5},\qquad
                  a_{n+1}=\frac{2a_n}{1-a_n^2}
                \]
                とする.
                \par\smallskip\noindent
                (1) \bookterm{general}{一般項}の候補を求め,すべての段階で分母が$\displaystyle 0$にならず,
                \bookterm{sequence}{数列}が定義できることを示せ.
                \par\smallskip\noindent
                (2) 周期の最小値を求めよ.
          % E13: ニュートン法型・初項が収束先より小さい場合
          \item \bookterm{sequence}{数列}$\displaystyle \{a_n\}$を
                \[
                  a_1=1,\qquad
                  a_{n+1}=\frac{a_n^2+3}{2a_n}
                \]
                によって定める.
                \par\smallskip\noindent
                (1) すべての項が正であることを示せ.
                \par\smallskip\noindent
                (2) $\displaystyle b_n=\frac{a_n-\sqrt3}{a_n+\sqrt3}$とおき,
                $\displaystyle b_n,a_n$の\bookterm{general}{一般項}を求めよ.
                \par\smallskip\noindent
                (3) $\displaystyle n\ge2$で$\displaystyle a_n>\sqrt3$かつ
                $\displaystyle a_{n+1}<a_n$であることを示し,\,$\displaystyle \lim_{n\to\infty}a_n$を求めよ.
          % E14: 初項依存型・不動点と2周期点
          \item 実数$\displaystyle c$を初項として,\bookterm{sequence}{数列}$\displaystyle \{a_n\}$を
                \[
                  a_1=c,\qquad a_{n+1}=a_n^2-1
                \]
                によって定める.
                \par\smallskip\noindent
                (1) \bookterm{sequence}{数列}が定\bookterm{sequence}{数列}となる$\displaystyle c$をすべて求めよ.
                \par\smallskip\noindent
                \BookMathLayout{(2) 定\bookterm{sequence}{数列}ではなく,すべての$\displaystyle n$で$\displaystyle a_{n+2}=a_n$を満たす
                $\displaystyle c$をすべて求め,それぞれの場合の\bookterm{general}{一般項}を表せ.}{(2) 定\bookterm{sequence}{数列}ではなく,\\すべての$\displaystyle n$で$\displaystyle a_{n+2}=a_n$を満たす
                $\displaystyle c$をすべて求め,それぞれの場合の\bookterm{general}{一般項}を表せ.}
          % E15: 初項依存型・有界になる初項の範囲
          \item 実数$\displaystyle c$を初項として,\bookterm{sequence}{数列}$\displaystyle \{a_n\}$を
                \[
                  a_1=c,\qquad a_{n+1}=a_n^2-2
                \]
                によって定める.
                \par\smallskip\noindent
                (1) $\displaystyle c=-2,-1,0,2$のそれぞれについて,最初の4項を求めよ.
                \par\smallskip\noindent
                (2) $\displaystyle -2\le c\le2$のとき,\,$\displaystyle c=2\cos\theta$
                （$\displaystyle 0\le\theta\le\pi$）とおき,\bookterm{general}{一般項}を求めよ.
                また,すべての項が$\displaystyle -2$以上$\displaystyle 2$以下であることを示せ.
                \par\smallskip\noindent
                (3) $\displaystyle |c|>2$のとき,\,$\displaystyle a_n\to+\infty$を示せ.
                以上から,\bookterm{sequence}{数列}が有界となる$\displaystyle c$の範囲を求めよ.
        \end{enumerate}

      \subsubsection{方針}\label{節：発展パターン・特殊性質型の演習：方針}
        \begin{enumerate}[label=\textbf{問\arabic*.},leftmargin=3.8em,format=\bookitemlink{practice-special}{hint}{problem}]
          % E1
          \item 整数を$\displaystyle 2$で割った商と余りに分け,
                $\displaystyle \left\lfloor\frac{\lfloor x\rfloor}{2}\right\rfloor
                =\left\lfloor\frac{x}{2}\right\rfloor$となることを確かめる.
                初めて$\displaystyle 0$になる位置は,\,$\displaystyle 2026$と$\displaystyle 2$の累乗を比べればよい.
          % E2
          \item \bookterm{difference}{階差}を足すと$\displaystyle \sum_{k=1}^{n}\lfloor\frac{k}{3}\rfloor$になる.
                同じ値を取る添字を三つずつまとめ,最後の群を$\displaystyle r+1$項として数える.
          % E3
          \item $\displaystyle x-5\lfloor\frac{x}{5}\rfloor$は整数$\displaystyle x$を$\displaystyle 5$で割った余りである.
                項の値が五つに絞れたら,それぞれが次にどの値へ移るかを表にする.
          % E4
          \item $\displaystyle d_{n+1}-d_n$を作ると元の\bookterm{recurrence}{漸化式}の左辺が現れる.
                その後,\,$\displaystyle e_n=a_{n+1}-a_n$を求め,もう一度足して$\displaystyle a_n$へ戻す.
          % E5
          \item $\displaystyle d_{n+1}-3d_n$を展開すると$\displaystyle 0$になる.
                得られた3項間の式を満たす$\displaystyle 3^{n-1}$を引き,残りを解く.
          % E6
          \item \bookterm{characteristic}{特性方程式}の根の絶対値と\bookterm{argument}{偏角}を調べる.
                最初の二つの項に合う$\displaystyle \sin$の式を作り,実際の値の並びから周期を確かめる.
          % E7
          \item \bookterm{root}{特性根}の絶対値は$\displaystyle 1$未満である.
                \bookterm{polar}{極形式}から\bookterm{trigonometric}{三角関数}の\bookterm{general}{一般項}を作り,その絶対値を\bookterm{geometric}{等比数列}で上から押さえる.
          % E8
          \item 正値を確かめてから逆数を取る.
                奇数番目と偶数番目を別々の\bookterm{sequence}{数列}にすると,二つの\bookterm{arithmetic}{等差数列}が得られる.
          % E9
          \item $\displaystyle a_n^2-2a_n=(a_n-1)^2-1$と\bookterm{square}{平方完成}する.
                $\displaystyle b_1=2+2^{-1}$と書けば,恒等式によって指数が毎回2倍になる.
          % E10
          \item 展開すると,交差項が$\displaystyle 3(u+u^{-1})$であることが分かる.
                初項を$\displaystyle 2+2^{-1}$と表し,指数を毎回3倍にする候補を作る.
          % E11
          \item 余弦の倍角公式で,角度を毎回2倍にする候補を作る.
                3回倍にした角度と元の角度との差が$\displaystyle 2\pi$の整数倍になることを調べる.
          % E12
          \item 正接の倍角公式を使う.
                候補の角度が$\displaystyle \frac{\pi}{2}$に$\displaystyle \pi$の整数倍を加えた値になるかを,
                整数の偶奇で判断する.次の角度でも正接が定義できれば,\bookterm{recurrence}{漸化式}の分母も$\displaystyle 0$にならない.
          % E13
          \item $\displaystyle a_{n+1}\pm\sqrt3$の分子は平方になる.比を取って$\displaystyle b_{n+1}=b_n^2$を導く.
                $\displaystyle b_1<0$なので,\bookterm{logarithm}{対数}を取らずにべき乗の形を帰納法で確かめる.
                減少を示す範囲が$\displaystyle n\ge2$であることにも注意する.
          % E14
          \item 定\bookterm{sequence}{数列}には$\displaystyle a_2=a_1$,周期2の候補には$\displaystyle a_3=a_1$が必要である.
                後者の方程式を因数分解し,定\bookterm{sequence}{数列}の初項を除く.残った初項が実際に二つの値を往復することも確かめる.
          % E15
          \item $\displaystyle |c|\le2$では$\displaystyle 2\cos\theta$という表示を使う.
                $\displaystyle |c|>2$では第2項が$\displaystyle 2$を超えるので,
                $\displaystyle a_{n+1}-2=(a_n-2)(a_n+2)$から下からの評価を作る.
        \end{enumerate}

      \subsubsection{解答}\label{節：発展パターン・特殊性質型の演習：解答}
        \begin{enumerate}[label=\textbf{問\arabic*.},leftmargin=3.8em,format=\bookitemlink{practice-special}{answer}{problem}]
          % E1
          \item (1)
                \[
                  a_n=\left\lfloor\frac{2026}{2^{n-1}}\right\rfloor.
                \]
                (2) $\displaystyle n=12$である.
          % E2
          \item
                \[
                  a_n=\frac{3q(q-1)}2+q(r+1),
                \]
                \[
                  n=3q+r,\qquad q\ge0,\quad r=0,1,2.
                \]
          % E3
          \item (1) 整数性と\bookterm{floor-function}{床関数}の定義から,各項は$\displaystyle 0$以上$\displaystyle 4$以下の整数となる.
                \par\smallskip\noindent
                (2) $\displaystyle k$を$\displaystyle 0$以上の整数として,
                \[
                  a_{4k+1}=1,\quad a_{4k+2}=3,
                \]
                \[
                  a_{4k+3}=2,\quad a_{4k+4}=0.
                \]
                周期の最小値は$\displaystyle 4$である.
          % E4
          \item (1) $\displaystyle d_n=6n+6$.
                \par\smallskip\noindent
                (2) $\displaystyle a_n=n^3+1$.
          % E5
          \item (1) $\displaystyle d_n=2\cdot3^{n-1}$.
                \par\smallskip\noindent
                (2)
                \[
                  a_n=1+2^{n-1}+3^{n-1}.
                \]
          % E6
          \item (1)
                \[
                  a_n=\frac2{\sqrt3}\sin\frac{(n-1)\pi}{3}.
                \]
                (2) 周期の最小値は$\displaystyle 6$である.
          % E7
          \item
                \[
                  a_n=\left(\frac1{\sqrt2}\right)^{n-1}
                    \left(\cos t_n-\sin t_n\right),
                \]
                \[
                  t_n=\frac{(n-1)\pi}{4},\qquad
                  \lim_{n\to\infty}a_n=0.
                \]
          % E8
          \item (1) すべての項は正であり,
                \[
                  a_{2k-1}=\frac1k,\qquad
                  a_{2k}=\frac2{2k-1}\quad(k\ge1).
                \]
                (2) \bookterm{limit}{極限}は$\displaystyle 0$である.
          % E9
          \item (1) $\displaystyle b_1=\frac52$,\,$\displaystyle b_{n+1}=b_n^2-2$.
                \par\smallskip\noindent
                (2)
                \[
                  a_n=1+2^{\,2^{n-1}}+2^{-2^{n-1}}.
                \]
          % E10
          \item (1) 左辺を展開すると右辺を得る.
                \par\smallskip\noindent
                (2)
                \[
                  a_n=2^{\,3^{n-1}}+2^{-3^{n-1}}.
                \]
          % E11
          \item (1)
                \[
                  a_n=\cos\left(2^{n-1}\frac{2\pi}{7}\right).
                \]
                (2) 周期の最小値は$\displaystyle 3$である.
          % E12
          \item (1)
                \[
                  a_n=\tan\left(2^{n-1}\frac{\pi}{5}\right).
                \]
                すべての項で正接が定義でき,\bookterm{recurrence}{漸化式}の分母も$\displaystyle 0$にならない.
                \par\smallskip\noindent
                (2) 周期の最小値は$\displaystyle 4$である.
          % E13
          \item (1) 初項が正で,正の項から次の正の項が定まる.
                \par\smallskip\noindent
                (2) $\displaystyle q=\sqrt3-2$とおくと,
                \[
                  b_n=q^{\,2^{n-1}},\qquad
                  a_n=\sqrt3\,\frac{1+q^{\,2^{n-1}}}{1-q^{\,2^{n-1}}}.
                \]
                (3) $\displaystyle n\ge2$で設問の二つの不等式が成り立ち,
                \[
                  \lim_{n\to\infty}a_n=\sqrt3.
                \]
          % E14
          \item (1)
                \[
                  c=\frac{1+\sqrt5}{2},\quad
                  \frac{1-\sqrt5}{2}.
                \]
                それぞれ$\displaystyle a_n=c$である.
                \par\smallskip\noindent
                (2) $\displaystyle c=0,-1$であり,どちらも
                \[
                  a_n=-\frac12+\left(c+\frac12\right)(-1)^{n-1}
                \]
                と表せる.周期の最小値は$\displaystyle 2$である.
          % E15
          \item (1)
                \[
                  \begin{array}{c|rrrr}
                    c&a_1&a_2&a_3&a_4\\ \hline
                    -2&-2&2&2&2\\
                    -1&-1&-1&-1&-1\\
                     0&0&-2&2&2\\
                     2&2&2&2&2
                  \end{array}
                \]
                (2)
                \[
                  a_n=2\cos(2^{n-1}\theta),\qquad -2\le a_n\le2.
                \]
                (3) $\displaystyle n\ge2$で
                \[
                  a_n\ge2+4^{n-2}(c^2-4)
                \]
                となり,\,$\displaystyle a_n\to+\infty$である.
                有界となる範囲は$\displaystyle -2\le c\le2$である.
        \end{enumerate}

      \subsubsection{解法}\label{節：発展パターン・特殊性質型の演習：解法}
        \begin{enumerate}[label=\textbf{問\arabic*.},leftmargin=3.8em,format=\bookitemlink{practice-special}{solution}{problem}]
          % E1
          \item \textbf{(1)~\bookterm{floor-function}{床関数}の中の整数部分を分ける.}
                \sectionref{節：ガウス型}の整数性に着目する.
                非負の実数$\displaystyle x$について,整数$\displaystyle \lfloor x\rfloor$を
                \[
                  \lfloor x\rfloor=2q+r,\qquad r=0,1
                \]
                と表す.\bookterm{floor-function}{床関数}の定義を使って$\displaystyle 2$で割れば,
                \begin{align*}
                  \frac x2&\ge q+\frac r2\ge q,\\
                  \frac x2&<q+\frac{r+1}{2}\le q+1.
                \end{align*}
                したがって,
                \[
                  \left\lfloor\frac{x}{2}\right\rfloor=q
                    =\left\lfloor\frac{\lfloor x\rfloor}{2}\right\rfloor.
                \]
                $\displaystyle a_1=\lfloor2026\rfloor$であり,この恒等式を繰り返し使うと,
                \[
                  a_n=\left\lfloor\frac{2026}{2^{n-1}}\right\rfloor
                \]
                が帰納法で得られる.

                \smallskip
                \textbf{(2)~切り捨てる前の値が1未満になる位置を見る.}
                $\displaystyle 2^{10}=1024\le2026<2048=2^{11}$だから,
                \[
                  a_{11}=1,\qquad a_{12}=0.
                \]
                $\displaystyle n\le11$では切り捨てる前の値が$\displaystyle 1$以上なので,
                初めて$\displaystyle 0$となるのは$\displaystyle n=12$である.
                毎回の切り捨てを一つの\bookterm{floor-function}{床関数}にまとめられる理由は,割る数が正の整数であることにある.
          % E2
          \item \textbf{同じ値を取る添字をまとめて足す.}
                \sectionref{節：ガウス型}の群\bookterm{sequence}{数列}の見方を使う.
                \bookterm{difference}{階差}を足すと,
                \[
                  a_n=\sum_{k=1}^{n}\left\lfloor\frac{k}{3}\right\rfloor.
                \]
                $\displaystyle k=1,2$では値は$\displaystyle 0$である.
                正の整数$\displaystyle j$について,\,$\displaystyle k=3j,3j+1,3j+2$の三つで値が$\displaystyle j$になる.
                $\displaystyle n=3q+r$とおくと,\,$\displaystyle q\ge1$の場合,
                値が$\displaystyle 1,2,\ldots,q-1$となる群がすべてそろい,
                値が$\displaystyle q$となる最後の群が$\displaystyle r+1$項だけある.
                よって,
                \begin{align*}
                  a_n&=3\sum_{j=1}^{q-1}j+q(r+1)\\
                     &=\frac{3q(q-1)}2+q(r+1).
                \end{align*}
                $\displaystyle q=0$,すなわち$\displaystyle n=0,1,2$でも,
                実際の値とこの式の値はともに$\displaystyle 0$である.
                最後の群を常に三つとして数えると誤るので,余り$\displaystyle r$で途中までの項数を表している.
          % E3
          \item \textbf{(1)~\bookterm{floor-function}{床関数}を余りとして読み替える.}
                \sectionref{節：ガウス型}では,\bookterm{floor}{ガウス記号}が整数問題の道具にもなる.
                整数$\displaystyle x$について\bookterm{floor-function}{床関数}の定義から,
                \[
                  0\le x-5\left\lfloor\frac{x}{5}\right\rfloor<5.
                \]
                この値は整数だから$\displaystyle 0,1,2,3,4$のいずれかである.
                初項は整数であり,整数の項から次の整数の項が定まるので,
                $\displaystyle x=2a_n+1$として帰納法を使えばよい.

                \smallskip
                \textbf{(2)~五つの値の更新を調べる.}
                (1)より,次の表だけを調べれば十分である.
                \[
                  \begin{array}{c|rrrrr}
                    a_n&0&1&2&3&4\\ \hline
                    a_{n+1}&1&3&0&2&4
                  \end{array}
                \]
                初項$\displaystyle 1$からは$\displaystyle 1,3,2,0,1,\ldots$と進む.
                どの値からも次の値は一つに決まるため,この四つの値の繰り返しが帰納法で続く.
                よって解答の\bookterm{general}{一般項}を得る.
                最初の四つの値がすべて異なるので,周期の最小値は$\displaystyle 4$である.
                \bookterm{sequence}{数列}全体の大きさを追う代わりに,取り得る値を先に絞ったことが解法の鍵である.
          % E4
          \item \textbf{(1)~4項間の関係を\bookterm{difference}{階差}の関係にする.}
                \sectionref{節：隣接4項間漸化式型}と同じく,\bookterm{difference}{階差}の組合せを見る.
                設問の$\displaystyle d_n$について,
                \begin{align*}
                  d_{n+1}-d_n
                    &=a_{n+3}-3a_{n+2}\\
                    &\quad+3a_{n+1}-a_n\\
                    &=6.
                \end{align*}
                $\displaystyle d_1=28-2\cdot9+2=12$だから,
                \[
                  d_n=12+6(n-1)=6n+6.
                \]

                \smallskip
                \textbf{(2)~\bookterm{difference}{階差}を順に足して元の\bookterm{sequence}{数列}へ戻す.}
                $\displaystyle e_n=a_{n+1}-a_n$とおくと,
                \[
                  e_{n+1}-e_n=d_n=6n+6,
                \]
                かつ$\displaystyle e_1=9-2=7$である.したがって,
                \begin{align*}
                  e_n&=7+\sum_{k=1}^{n-1}(6k+6)\\
                     &=3n^2+3n+1.
                \end{align*}
                もう一度足せば,
                \begin{align*}
                  a_n&=2+\sum_{k=1}^{n-1}(3k^2+3k+1)\\
                     &=n^3+1.
                \end{align*}
                ここでは$\displaystyle n=1$の空の和を$\displaystyle 0$とする.
                三つの初期値を,\,$\displaystyle d_1,e_1,a_1$の順に使っている.
                \bookterm{difference}{階差}を取ると次数が下がり,足し合わせると元の\bookterm{sequence}{数列}へ戻る関係が見える.
          % E5
          \item \textbf{(1)~等比型になる3項の組合せを作る.}
                \sectionref{節：隣接4項間漸化式型}の方法を使う.
                \bookterm{characteristic}{特性方程式}の左辺は
                \[
                  x^3-6x^2+11x-6=(x-1)(x-2)(x-3)
                \]
                であり,\,$\displaystyle 1,2$に対応する二次式から設問の$\displaystyle d_n$が作られる.
                実際に展開すると,
                \[
                  d_{n+1}-3d_n
                    =a_{n+3}-6a_{n+2}+11a_{n+1}-6a_n=0.
                \]
                $\displaystyle d_1=14-3\cdot6+2\cdot3=2$より,
                \[
                  d_n=2\cdot3^{n-1}.
                \]
                初期値が$\displaystyle 0$ではないので,この右辺を残して次の変形を行う.

                \smallskip
                \textbf{(2)~右辺を作る\bookterm{sequence}{数列}を引く.}
                $\displaystyle 3^{n-1}$に同じ組合せを施すと,
                \[
                  3^{n+1}-3\cdot3^n+2\cdot3^{n-1}=2\cdot3^{n-1}.
                \]
                そこで$\displaystyle b_n=a_n-3^{n-1}$とおけば,
                \[
                  b_{n+2}-3b_{n+1}+2b_n=0.
                \]
                \bookterm{root}{特性根}は$\displaystyle 1,2$だから$\displaystyle b_n=A+B2^{n-1}$とおける.
                $\displaystyle b_1=2,b_2=3$から$\displaystyle A=B=1$となり,
                \[
                  a_n=1+2^{n-1}+3^{n-1}
                \]
                を得る.4項間を3項間へ,さらに既習の解法へ順に帰着させた.
          % E6
          \item \textbf{(1)~虚数の\bookterm{root}{特性根}を角度で表す.}
                \sectionref{節：虚数解型隣接3項間漸化式}を使う.
                \bookterm{characteristic}{特性方程式}$\displaystyle x^2-x+1=0$の根は,
                \[
                  \frac{1\pm\sqrt3\,i}{2}
                   =\cos\frac\pi3\pm i\sin\frac\pi3.
                \]
                絶対値が$\displaystyle 1$なので,
                \[
                  a_n=C\cos\frac{(n-1)\pi}{3}
                      +D\sin\frac{(n-1)\pi}{3}
                \]
                とおける.\,$\displaystyle a_1=0$より$\displaystyle C=0$であり,
                $\displaystyle a_2=1$より$\displaystyle D=\frac2{\sqrt3}$である.
                得られた\bookterm{general}{一般項}が\bookterm{recurrence}{漸化式}を満たすことは,\,$\displaystyle \sin$の\bookterm{addition}{加法定理}からも確認できる.

                \smallskip
                \textbf{(2)~周期の最小値まで確かめる.}
                \bookterm{general}{一般項}より$\displaystyle a_{n+6}=a_n$であり,最初の6項は
                \[
                  0,1,1,0,-1,-1
                \]
                である.周期が$\displaystyle 1,2,4,5$なら$\displaystyle a_{1+p}=a_1$が成り立たない.
                周期が$\displaystyle 3$なら$\displaystyle a_5=a_2$が必要だが,
                $\displaystyle -1\ne1$である.
                したがって周期の最小値は$\displaystyle 6$である.

                \smallskip
                \noindent \textbf{【別解】}\\
                \bookterm{recurrence}{漸化式}を一つ進めると,
                \[
                  a_{n+3}=a_{n+2}-a_{n+1}=-a_n.
                \]
                さらに3項進めれば$\displaystyle a_{n+6}=a_n$となる.
                \bookterm{root}{特性根}の角度と,直接の式変形による周期の証明が対応している.
          % E7
          \item \textbf{\bookterm{general}{一般項}の振幅を調べて\bookterm{limit}{極限}を求める.}
                \sectionref{節：虚数解型隣接3項間漸化式}では,\bookterm{root}{特性根}の絶対値も見る.
                \bookterm{characteristic}{特性方程式}$\displaystyle x^2-x+\frac12=0$の根は,
                \[
                  \frac{1\pm i}{2}
                    =\frac1{\sqrt2}
                      \left(\cos\frac\pi4\pm i\sin\frac\pi4\right).
                \]
                したがって$\displaystyle t_n=\frac{(n-1)\pi}{4}$として,
                \[
                  a_n=\left(\frac1{\sqrt2}\right)^{n-1}
                    (C\cos t_n+D\sin t_n)
                \]
                とおける.\,$\displaystyle a_1=1$から$\displaystyle C=1$であり,
                $\displaystyle a_2=0$から$\displaystyle D=-1$である.
                よって解答の\bookterm{general}{一般項}を得る.\bookterm{trigonometric}{三角関数}の絶対値は$\displaystyle 1$以下なので,
                \[
                  |a_n|\le2\left(\frac1{\sqrt2}\right)^{n-1}\longrightarrow0.
                \]
                したがって$\displaystyle a_n\to0$となる.
                角度は一定ずつ進み,振幅は\bookterm{geometric}{等比数列}として縮んでいる.
          % E8
          \item \textbf{(1)~正値を確かめ,奇数番目と偶数番目を分ける.}
                \sectionref{節：非隣接2項間漸化式型}の着目点は,添字が一つ飛んでいることである.
                $\displaystyle a_1,a_2>0$であり,\,$\displaystyle a_n>0$なら
                $\displaystyle \frac{a_n}{1+a_n}>0$となる.
                奇数番目と偶数番目のそれぞれで帰納法を使えば,すべての項が正である.
                よって逆数を取ることができ,
                \[
                  \frac1{a_{n+2}}=\frac1{a_n}+1.
                \]
                $\displaystyle b_n=\frac1{a_n}$とおくと,
                \[
                  b_1=1,\quad b_2=\frac12,\quad b_{n+2}=b_n+1.
                \]
                したがって,
                \[
                  b_{2k-1}=k,\qquad b_{2k}=k-\frac12.
                \]
                逆数を取って解答の\bookterm{general}{一般項}を得る.

                \smallskip
                \textbf{(2)~二つの部分列を合わせて\bookterm{limit}{極限}を調べる.}
                奇数番目と偶数番目のどちらも$\displaystyle 0$へ近づく.
                全体としても,\,$\displaystyle n\ge2$では\bookterm{general}{一般項}から
                \[
                  0<a_n\le\frac2{n-1}
                \]
                となるので,\,$\displaystyle a_n\to0$である.
                添字を分けたあとには,各部分列の答えを元の添字へ戻して確認する.
          % E9
          \item \textbf{(1)~基準となる値をずらして構造を出す.}
                \bookterm{square}{平方完成}して$\displaystyle b_n=a_n-1$とおくと,
                \begin{align*}
                  b_{n+1}&=a_n^2-2a_n-1\\
                         &=(a_n-1)^2-2=b_n^2-2.
                \end{align*}
                初項は$\displaystyle b_1=\frac52$である.

                \smallskip
                \textbf{(2)~基本対称式の指数だけを変える.}
                \sectionref{節：基本対称式型}と同じ形が現れた.
                $\displaystyle \frac52=2+2^{-1}$であるから,
                \[
                  B_n=2^{\,2^{n-1}}+2^{-2^{n-1}}
                \]
                を候補とする.設問の恒等式に$\displaystyle u=2^{\,2^{n-1}}$を入れると,
                \[
                  B_{n+1}=B_n^2-2.
                \]
                $\displaystyle B_1=b_1$でもあるので,帰納法により$\displaystyle b_n=B_n$となる.
                よって$\displaystyle a_n=1+B_n$である.
                元の式で\bookterm{logarithm}{対数}を取る前に\bookterm{square}{平方完成}したことで,対称式の恒等式に一致する形が見つかった.
          % E10
          \item \textbf{(1)~交差項が打ち消されることを確認する.}
                \sectionref{節：基本対称式型}の三次の恒等式を使う.
                展開すると,
                \[
                  (u+u^{-1})^3=u^3+3u+3u^{-1}+u^{-3}
                \]
                であるから,\,$\displaystyle 3(u+u^{-1})$を引けばよい.

                \smallskip
                \textbf{(2)~指数の3倍を繰り返す.}
                $\displaystyle a_1=2+2^{-1}$に合う候補を
                \[
                  A_n=2^{\,3^{n-1}}+2^{-3^{n-1}}
                \]
                とする.(1)に$\displaystyle u=2^{\,3^{n-1}}$を代入すると,
                \[
                  A_{n+1}=A_n^3-3A_n
                \]
                となる.初項と\bookterm{recurrence}{漸化式}が一致するので,帰納法より$\displaystyle a_n=A_n$である.
                $\displaystyle -3a_n$の項が交差項を正確に打ち消すため,同じ対称式の形を保って次へ進める.
          % E11
          \item \textbf{(1)~角度の倍増を使って候補を作る.}
                \sectionref{節：三角関数型}の余弦の倍角公式を使う.
                $\displaystyle \theta_n=2^{n-1}\frac{2\pi}{7}$とおき,
                $\displaystyle A_n=\cos\theta_n$とすれば,
                \[
                  A_{n+1}=\cos(2\theta_n)=2A_n^2-1.
                \]
                初項も一致するので,帰納法により$\displaystyle a_n=A_n$である.
                これは角度の選び方によって解を作ったのであり,
                余弦の値から角度が一意に定まると主張しているわけではない.

                \smallskip
                \textbf{(2)~一周の整数倍だけ違う角度を比較する.}
                3回進めると,
                \[
                  \theta_{n+3}-\theta_n=2^n\pi
                      =2\pi\cdot2^{n-1}.
                \]
                これは$\displaystyle 2\pi$の整数倍だから,\,$\displaystyle a_{n+3}=a_n$である.
                最初の三つの値は,
                \[
                  \cos\frac{2\pi}{7},\quad
                  \cos\frac{4\pi}{7},\quad
                  \cos\frac{6\pi}{7}
                \]
                となる.\,$\displaystyle \cos$は$\displaystyle 0\le\theta\le\pi$で狭義に減少するので,
                三つの値は異なり,周期の最小値は$\displaystyle 3$である.
          % E12
          \item \textbf{(1)~正接と\bookterm{recurrence}{漸化式}の両方の分母を確認する.}
                \sectionref{節：三角関数型}の正接の倍角公式から,
                \[
                  \theta_n=2^{n-1}\frac\pi5,\qquad A_n=\tan\theta_n
                \]
                を候補とする.\,$\displaystyle \cos\theta_n=0$となるには,
                ある整数$\displaystyle j$について
                \[
                  2^{n-1}\frac\pi5=\frac\pi2+j\pi,
                  \quad\text{すなわち}\quad 2^n=5(2j+1)
                \]
                が必要である.左辺は偶数,右辺は奇数なので,これは起こらない.
                よって$\displaystyle A_n$はすべての$\displaystyle n\ge1$で定義できる.
                さらに,
                \[
                  1-A_n^2=\frac{\cos(2\theta_n)}{\cos^2\theta_n}
                           =\frac{\cos\theta_{n+1}}{\cos^2\theta_n}\ne0.
                \]
                したがって,倍角公式を各段階で使うことができ,
                \[
                  A_{n+1}=\frac{2A_n}{1-A_n^2}
                \]
                となる.初項も一致するから,帰納法で$\displaystyle a_n=A_n$を得る.

                \smallskip
                \textbf{(2)~正接の周期を角度の差へ戻す.}
                4回進めると,
                \[
                  \theta_{n+4}-\theta_n=3\cdot2^{n-1}\pi.
                \]
                $\displaystyle \tan$の周期は$\displaystyle \pi$なので,\,$\displaystyle a_{n+4}=a_n$である.
                最初の四つの値は,
                \[
                  \tan\frac\pi5,\quad \tan\frac{2\pi}{5},\quad
                  -\tan\frac\pi5,\quad -\tan\frac{2\pi}{5}
                \]
                である.\,$\displaystyle 0<\frac\pi5<\frac{2\pi}{5}<\frac\pi2$で正接は狭義に増加するから,
                この四つは異なる.よって周期の最小値は$\displaystyle 4$である.
          % E13
          \item \textbf{(1)~変形する前に正値を確かめる.}
                $\displaystyle a_1=1>0$であり,\,$\displaystyle a_n>0$なら分子・分母がともに正である.
                よって帰納法によりすべての項が正であり,\,$\displaystyle a_n+\sqrt3$も$\displaystyle 0$にならない.

                \smallskip
                \textbf{(2)~二つの平方の比を作る.}
                \sectionref{節：ニュートン法型}の$\displaystyle \sqrt2$を$\displaystyle \sqrt3$に置き換えた形である.
                \[
                  a_{n+1}-\sqrt3=\frac{(a_n-\sqrt3)^2}{2a_n},
                \]
                \[
                  a_{n+1}+\sqrt3=\frac{(a_n+\sqrt3)^2}{2a_n}.
                \]
                共通の分母を消すために比を取ると,
                \[
                  b_{n+1}=b_n^2.
                \]
                初項は
                \[
                  b_1=\frac{1-\sqrt3}{1+\sqrt3}=\sqrt3-2=q.
                \]
                $\displaystyle q<0$なので実数の\bookterm{logarithm}{対数}を最初から取ることはできないが,
                繰り返し二乗する関係から,帰納法により
                \[
                  b_n=q^{\,2^{n-1}}
                \]
                が得られる.\,$\displaystyle b_n$の定義を$\displaystyle a_n$について解けば,
                \[
                  a_n=\sqrt3\,\frac{1+b_n}{1-b_n}.
                \]
                $\displaystyle -1<q<0$なので$\displaystyle |b_n|<1$であり,
                分母$\displaystyle 1-b_n$は常に正である.

                \smallskip
                \textbf{(3)~第2項からの減少を調べる.}
                $\displaystyle n\ge2$では$\displaystyle b_n>0$なので,\bookterm{general}{一般項}から$\displaystyle a_n>\sqrt3$である.
                また,
                \[
                  a_{n+1}-a_n=\frac{3-a_n^2}{2a_n}<0
                  \quad(n\ge2).
                \]
                $\displaystyle |q|<1$から$\displaystyle b_n\to0$であり,\,$\displaystyle a_n\to\sqrt3$となる.
                初項は収束先より小さいが,第2項でその上へ移り,そこから減少して近づく.
                誤差の式$\displaystyle a_{n+1}-\sqrt3=\frac{(a_n-\sqrt3)^2}{2a_n}$も,
                第2項から上側に移ることと,誤差が急速に小さくなることを説明している.
          % E14
          \item \textbf{(1)~定\bookterm{sequence}{数列}を\bookterm{fixed}{不動点}として求める.}
                \sectionref{節：初項依存型}の考え方を使う.
                定\bookterm{sequence}{数列}には$\displaystyle c=c^2-1$が必要だから,
                \[
                  c^2-c-1=0,\qquad c=\frac{1\pm\sqrt5}{2}.
                \]
                これらの初項なら毎回同じ値になるので,十分条件でもある.

                \smallskip
                \textbf{(2)~2回進めて戻る初項を探す.}
                $\displaystyle a_3=a_1$から,
                \begin{align*}
                  (c^2-1)^2-1&=c,\\
                  c^4-2c^2-c&=0,\\
                  c(c+1)(c^2-c-1)&=0.
                \end{align*}
                (1)の定\bookterm{sequence}{数列}を除くと$\displaystyle c=0,-1$だけが残る.
                実際,更新は$\displaystyle 0\mapsto-1\mapsto0$となるので,
                どちらの初項もこの二つの値を交互に取る.
                その\bookterm{general}{一般項}は
                \[
                  a_n=-\frac12+\left(c+\frac12\right)(-1)^{n-1}
                  \quad(c=0,-1)
                \]
                と表せる.二つの値が異なるから周期の最小値は$\displaystyle 2$である.
                この方程式から得た四つの初項について特別な振る舞いを示したのであり,
                他の初項で同じ\bookterm{general}{一般項}が使えるとは結論していない.
          % E15
          \item \textbf{(1)~初項を変えて同じ規則の振る舞いを見る.}
                各初項を代入すると解答の表を得る.
                $\displaystyle c=-1,2$では最初から定\bookterm{sequence}{数列}となり,
                $\displaystyle c=-2,0$では数項進めたあとに$\displaystyle 2$で一定になる.
                同じ\bookterm{recurrence}{漸化式}でも,初項が違うと一定になるまでの経過が変わる.

                \smallskip
                \textbf{(2)~初項の範囲に合う表示を選ぶ.}
                \sectionref{節：初項依存型}から\sectionref{節：三角関数型}へつながる例である.
                $\displaystyle -2\le c\le2$なら$\displaystyle c=2\cos\theta$
                （$\displaystyle 0\le\theta\le\pi$）と表せる.
                $\displaystyle A_n=2\cos(2^{n-1}\theta)$とおくと,
                \begin{align*}
                  A_n^2-2
                    &=4\cos^2(2^{n-1}\theta)-2\\
                    &=2\cos(2^n\theta)=A_{n+1}.
                \end{align*}
                初項も一致するので,帰納法により$\displaystyle a_n=A_n$である.
                $\displaystyle |\cos x|\le1$から,すべての項で$\displaystyle |a_n|\le2$となる.

                \smallskip
                \textbf{(3)~範囲の外では2からの距離を下から押さえる.}
                $\displaystyle |c|>2$なら$\displaystyle a_2=c^2-2>2$である.
                $\displaystyle a_n>2$のとき,
                \[
                  a_{n+1}-2=(a_n-2)(a_n+2)>4(a_n-2)>0.
                \]
                よって帰納法で$\displaystyle a_n>2$が$\displaystyle n\ge2$のすべてで続く.
                この不等式を繰り返すと,
                \[
                  a_n-2\ge4^{n-2}(a_2-2)
                         =4^{n-2}(c^2-4).
                \]
                右辺は$\displaystyle +\infty$へ発散するので,\,$\displaystyle a_n\to+\infty$である.
                したがって有界となる初項はちょうど$\displaystyle -2\le c\le2$である.
                \bookterm{general}{一般項}が表せる範囲と,不等式で振る舞いを調べる範囲を,初項に応じて選んだ.
        \end{enumerate}
